%%%%%%%%%%  scan idx 434  |  folio 421  %%%%%%%%%%
base fidèlement plat quasi-compact. Si $i'$ est régulière, alors l'Idéal $J$ qui définit $i$ est de type fini, et $J/J^2$ est localement libre de type fini, car ces propriétés se descendent par morphismes fidèlement plats quasi-compacts. Si $f_1,\ldots,f_d$ est une suite de sections de $J$ relevant une base de $J/J^2$, le complexe de Koszul associé aux images inverses des $f_i$ est acyclique en degrés $>0$ d'après 1.4; comme c'est l'image inverse du complexe de Koszul associé aux $f_i$, ce dernier est acyclique en degrés $>0$, et $i$ est régulière.

\textbf{Définition 1.6.} Soit $i:Y\longrightarrow X$ une immersion définie dans un ouvert $U$ de $X$ par un Idéal $J$ de $\cO_X$; le $\cO_Y$-Module $J/J^2$ (indépendant de $U$) est appelé faisceau conormal de $Y$ dans $X$, et noté $N_{X/Y}$. Si $i$ est régulière, $J/J^2$ est localement libre de type fini, et son rang s'appelle codimension de $Y$ dans $X$, ou encore codimension de $i$.

Si on a deux immersions fermées : $Z\xrightarrow{j}Y\xrightarrow{i}X$, les faisceaux conormaux forment une suite exacte:
\begin{equation}\tag{1.6.1}
 j^*(N_{X/Y})\longrightarrow N_{X/Z}\longrightarrow N_{Y/Z}\longrightarrow0.
\end{equation}
Soient en effet $I$ l'idéal de $Y$ dans $X$, $J$ l'idéal de $Z$ dans $X$, $K$ l'idéal de $Z$ dans $Y$. On a la suite exacte:
\[
0\longrightarrow I\longrightarrow J\longrightarrow K\longrightarrow0;
\]
d'où par tensorisation par $\cO_X/I=\cO_Y$, la suite exacte:
\[
 I/I^2\longrightarrow J\otimes_{\cO_X}\cO_Y\longrightarrow K\otimes_{\cO_X}\cO_Y\longrightarrow0;
\]
en tensorisant par $\cO_Z=\cO_X/J=\cO_Y/K$, on obtient la suite exacte:
\[
 j^*(I/I^2)\longrightarrow J\otimes_{\cO_X}\cO_Z\longrightarrow K\otimes_{\cO_Y}\cO_Z\longrightarrow0,
\]
qui est la suite cherchée.

